% Diagrama: Dos tipos de salida de Apimetro
% Columna izquierda: JSON Descriptivo (atributos alfanuméricos)
% Columna derecha: GeoJSON Geográfico (geometría + métricas)
% Usado en: 4-Capas-Codigo/4-1-Apimetro/4-1-4-Salidas.tex (fig:apimetro-salidas-tipos)
\begin{tikzpicture}[
  cabecera/.style={
    rectangle, draw=unamAzul, fill=unamAzul, text=white,
    minimum width=5.4cm, minimum height=0.9cm,
    align=center, font=\small\bfseries
  },
  cuerpo/.style={
    rectangle, draw=unamAzul!40, fill=unamAzul!8,
    minimum width=5.4cm, text width=5.0cm,
    align=left, inner sep=8pt, font=\footnotesize
  },
  uso/.style={
    rectangle, draw=unamOro!70, fill=unamOro!18,
    minimum width=5.4cm, minimum height=0.75cm,
    align=center, font=\footnotesize\itshape, text=unamAzul
  }
]

%% ── Columna izquierda: JSON Descriptivo ──────────────────────────
\node[cabecera] (jh) at (0, 0) {JSON Descriptivo};

\node[cuerpo, anchor=north] (jb) at (jh.south) {%
  \textbf{¿Qué contiene?}\\[4pt]
  \texttt{linea\_id} · \texttt{nombre}\\
  \texttt{sistema} · \texttt{color\_esp}\\
  \texttt{clasificacion}\\[4pt]
  \textit{Sin geometría}
};

\node[uso, anchor=north] (ju) at (jb.south)
  {tablas · fichas · simbología};

%% ── Columna derecha: GeoJSON Geográfico ──────────────────────────
\node[cabecera] (gh) at (6.6, 0) {GeoJSON Geográfico};

\node[cuerpo, anchor=north] (gb) at (gh.south) {%
  \textbf{¿Qué contiene?}\\[4pt]
  \texttt{coordinates} (Point / MultiLineString)\\
  \texttt{properties}\\
  \quad velocidad · frecuencia · capacidad\\[4pt]
  \textit{Listo para web, sin conversión}
};

\node[uso, anchor=north] (gu) at (gb.south)
  {mapas · ruteo · \textit{VFTModel}};

\end{tikzpicture}
